\appendix

\chapter{Formal Axiomatic Reference}
\label{app:axioms}

\begin{introduction}[Structure of the Formal System]
  \item Two \textbf{fundamental axioms}: Poisson causality
        (finite DAG + 4D sprinkling) and planarity preservation.
  \item Seven \textbf{derived theorems}: results that the original nine axioms
        reduced to postulates, now proven from
        the two axioms.
  \item Eight \textbf{laws}: operational rules governing the dynamics
        of the graph.
  \item Seven \textbf{theorems}: results mathematically derivable from
        the axioms and the laws.
  \item Ten \textbf{conjectures and predictions}: open problems that the
        software must still validate numerically or that require a
        pending analytical proof.
\end{introduction}

\noindent
The axioms, laws, and theorems constitute the formal system. The
simulation \texttt{prism\_simulation} implements and verifies them
numerically, but the mathematics is independent of the code.
The source-code references accompanying each entry serve
as a correspondence guide, not as formal justification.

% ═══════════════════════════════════════════════════════════════════════════════
\section{Fundamental Axioms}
% ═══════════════════════════════════════════════════════════════════════════════

\begin{axiom}[Poisson Causality]
\label{ax:causal-finitude}
The universe is a locally finite
partially ordered set $\mathcal{C} = (V, \prec)$---a Directed
Acyclic Graph (DAG), instantiated as a Poisson
sprinkling at density $\rho = 1$ over a Lorentzian manifold of
dimensionality $d = 4$. The set $V$ is finite, the
relation~$\prec$ is inherited from the light cone, and all
internal accumulation uses integer arithmetic
(\textbf{strict finitism}). Lorentz invariance is the
only symmetry compatible with the Poisson
sprinkling~\cite{bombelli2009}.

\medskip\noindent
\texttt{diamond.rs}: \texttt{sprinkle(n, rng)} generates exactly
$n$ points via a Poisson process. \texttt{is\_causal(a, b)}
implements the light-cone cut: $t_v - t_u > 0$ and
$(t_v - t_u)^2 > |\Delta\vec{x}|^2$.
\texttt{spectral.rs}: all internal accumulation uses
\texttt{u64}; the only \texttt{f64} division occurs at
the end: $P(t) = \text{count} / W$.
\end{axiom}

\begin{axiom}[Planarity Preservation]
\label{ax:planarity-preservation-app}
The Hasse diagram of~$\mathcal{C}$ preserves local
planarity: every emergent subdivision of~$K_5$ or~$K_{3,3}$ is
resolved by vertex absorption---the external node that
generates the obstruction contracts into the highest-degree pole of
the threatened Prism.

\medskip\noindent
See Axiom~\ref{ax:planarity-preservation} and
Theorem~\ref{thm:k2n-uniqueness} (Chapter~1).
\end{axiom}

% ═══════════════════════════════════════════════════════════════════════════════
\section{Derived Theorems}
\label{sec:derived-theorems}
% ═══════════════════════════════════════════════════════════════════════════════

\noindent
The following results were axioms in earlier versions of the
formal system. They are now proven from the two fundamental
axioms.

\begin{theorem}[Existence and Uniqueness of the Hasse Diagram (ex-Axiom)]
\label{thm:covering-differential}
Every finite DAG possesses a unique transitive reduction: the
\textbf{Hasse diagram}. The covering link $\covers$ is
the fundamental physical degree of freedom. Every edge $(u,v)$ such
that a path $u \to w \to v$ exists is redundant and is eliminated.
The Hasse diagram is the canonical description of
spacetime.
\end{theorem}

\begin{proof}[Sketch]
The transitive reduction is a unique operation on any
finite DAG (Axiom~\ref{ax:causal-finitude}). The
Stone--Kuratowski duality (Theorem~\ref{thm:stone-kuratowski}) shows that
$\covers$ is the atomic implication in the Heyting
algebra, and the transitive reduction is cut
elimination~\cite{gentzen1935}.
\end{proof}

\medskip\noindent
\texttt{diamond.rs}: \texttt{build\_hasse\_sparse()} eliminates
edges with $A^2[u,v] > 0$. \texttt{build\_hasse\_direct()} uses
incremental transitive reduction.

\begin{theorem}[Valence Saturation (ex-Axiom)]
\label{thm:valence-saturation}
The total degree of any event in the Hasse diagram is
bounded:
\begin{equation}
  \deg(u) \leq \mathcal{V}_{\max} = 15.
  \label{eq:ax-valence}
\end{equation}
\end{theorem}

\begin{proof}[Sketch]
Axiom~\ref{ax:causal-finitude} fixes the Poisson sprinkling
at $d = 4$, whence $\langle k\rangle = 2d = 8$. The
exact Poisson($8$) tail gives $\Pr(k > 15) \approx 0{.}8\,\%$.
See Theorem~\ref{thm:valence-limit} (Chapter~1).
\end{proof}

\medskip\noindent
\texttt{diamond.rs}: \texttt{const MAX\_HASSE\_DEGREE: usize = 15;}

\begin{theorem}[Causal Locality (ex-Axiom)]
\label{thm:causal-locality}
The search for Hasse links and the detection of
topological defects operate within a local horizon of
depth:
\begin{equation}
  d_{\max} = 4 \;\;\text{Hasse strata.}
  \label{eq:ax-locality}
\end{equation}
\end{theorem}

\begin{proof}[Sketch]
Prism diameter $= 2$ plus one hop in each direction for
$K_5$ threat detection: $d_{\max} = 2 + 2 = 4$.
See Theorem~\ref{thm:hasse-depth} (Chapter~1).
\end{proof}

\medskip\noindent
\texttt{diamond.rs}: \texttt{const MAX\_CAUSAL\_DEPTH: i16 = 4;}

\begin{theorem}[Bipartite Matter (ex-Axiom)]
\label{thm:bipartite-matter}
Matter is a complete bipartite subgraph $K_{2,N}$
called a \textbf{Causal Prism}, composed of:
\begin{enumerate}
  \item Two \textbf{poles} (origin $u$ and destination $v$) at Hasse
    distance $\geq 2$.
  \item $N \geq 3$ \textbf{intermediate} nodes $\{w_1,\ldots,w_N\}$,
    mutually disconnected (guaranteed by the
    triangle-free property), such that $u \covers w_i$ and
    $w_i \covers v$ for all $i$.
\end{enumerate}
The topological mass is the integer $M = N = |\text{intermediates}|$.
\end{theorem}

\begin{proof}[Sketch]
Theorem~\ref{thm:k2n-uniqueness} shows that $K_{2,N}$ is the
only non-trivial planar defect: triangle-freeness eliminates
$K_4$, planarity eliminates $K_{3,3}$, and $K_{1,N}$ is trivial.
\end{proof}

\medskip\noindent
\texttt{skyrmion.rs}: \texttt{struct CausalPrism \{origin,
destination, intermediates: Vec<usize>\}}.
\texttt{const MIN\_PRISM\_SHARED: usize = 3;}

\stoneref{The Prism is \emph{Modus Ponens} with $N$ intermediate lemmas
(Theorem~\ref{thm:antimatter-tollens}).}

\begin{theorem}[Kuratowski Threat (ex-Axiom)]
\label{thm:kuratowski-threat}
An external node $t \notin \mathcal{P}$ satisfying:
\begin{equation}
  t \sim u,\quad t \sim v,\quad |\{w_i \in W : t \sim w_i\}|
  \geq K_{\text{threat}} = 2,
  \label{eq:ax-threat}
\end{equation}
constitutes a $K_5$ threat and is contracted by absorption into the
highest-degree pole. This operation preserves the planarity
of the graph and is the discrete analogue of strong confinement.
\end{theorem}

\begin{proof}[Sketch]
Case analysis of Kuratowski's Theorem~\cite{kuratowski1930}:
$K_{\text{threat}} = 1$ is insufficient (4 vertices),
$K_{\text{threat}} = 3$ already violates planarity ($K_{3,3}$).
See Theorem~\ref{thm:kuratowski-threshold} (Chapter~1).
\end{proof}

\medskip\noindent
\texttt{skyrmion.rs}: \texttt{const PRISM\_THREAT: usize = 2;}

\begin{theorem}[Emergent Lorentz Invariance (ex-Axiom)]
\label{thm:lorentzian-causality}
The causal relation $u \prec v$ of the DAG is a pure logical arrow
($A \Rightarrow B$). A macroscopic observer measuring that
implication through continuous coordinates $(t, \vec{x})$ recovers
the Minkowski metric:
\begin{equation}
  u \prec v
  \quad\Longleftrightarrow\quad
  t_v - t_u > 0
  \;\;\text{and}\;\;
  (t_v - t_u)^2 > |\Delta\vec{x}|^2.
  \label{eq:ax-lorentz}
\end{equation}
Lorentz invariance is not imposed as an additional symmetry;
it is the only symmetry compatible with the Poisson sprinkling
declared in Axiom~\ref{ax:causal-finitude}.
\end{theorem}

\begin{proof}[Sketch]
Axiom~\ref{ax:causal-finitude} instantiates the DAG as a
Poisson sprinkling over a 4D Lorentzian manifold.
The Bombelli--Henson--Sorkin theorem~\cite{bombelli2009}
(Theorem~\ref{thm:lorentz-sprinkling}) shows that the
Poisson sprinkling is the only discretization
compatible with Lorentz invariance. Therefore, the
relation $\prec$ inherited from the light cone---which at the discrete level
is a logical arrow of the DAG---manifests macroscopically
as condition~\eqref{eq:ax-lorentz}.
\end{proof}

\medskip\noindent
\texttt{diamond.rs}: \texttt{is\_causal(a, b)} literally
implements~\eqref{eq:ax-lorentz}.

\begin{theorem}[Planck Density (ex-Axiom)]
\label{thm:planck-density}
The Poisson \emph{sprinkling} operates at the fundamental
density $\rho = 1$ event per Planck 4-volume:
\begin{equation}
  T = \Bigl(\frac{24\,N}{\pi}\Bigr)^{\!1/4},
  \qquad
  V = \frac{\pi\,T^4}{24},
  \qquad
  \rho = \frac{N}{V} = 1.
  \label{eq:ax-planck}
\end{equation}
\end{theorem}

\begin{proof}[Sketch]
Axiom~\ref{ax:causal-finitude} declares $\rho = 1$ in Planck
units. This is a choice of natural units (dimensional
analysis), not an independent degree of freedom.
\end{proof}

\medskip\noindent
\texttt{diamond.rs}: \texttt{let big\_t =
(24.0 * n as f64 / PI).powf(0.25);}

% ═══════════════════════════════════════════════════════════════════════════════
\section{Laws}
% ═══════════════════════════════════════════════════════════════════════════════

\subsection*{L1. Law of Transitive Reduction}
\phantomsection\label{law:transitive-reduction}

The edge $(u,v)$ is a Hasse link if and only if
$A^2[u,v] = 0$: no two-hop path exists. This
operation produces a \textbf{triangle-free} DAG by
mathematical necessity. Consequence: $K_4$ cliques cannot
exist---and the bipartite structure $K_{2,N}$ is the
maximum allowed.

\medskip\noindent
\texttt{diamond.rs}: \texttt{build\_hasse\_sparse()} computes
$A^2$ with sparse matrix multiplication (\texttt{sprs}) and
eliminates every edge with $A^2[i,j] > 0$.
\texttt{build\_hasse\_direct()} implements the same logic
incrementally: ``edge $(u,v)$ is a link iff no existing
link-child $z$ of $u$ satisfies $z \prec v$.''

\stoneref{The transitive reduction is cut elimination
(Theorem~\ref{thm:stone-kuratowski}) and Non-Contradiction
(Theorem~\ref{thm:non-contradiction}).}

\subsection*{L2. Law of Topological Mass}
\phantomsection\label{law:topological-mass}

The mass of a Causal Prism is the integer number of intermediate
nodes:
\begin{equation}
  M = N = |W| = |\{w : u \covers w \covers v\}|
  \in \mathbb{N},\quad N \geq 3.
  \label{eq:law-mass}
\end{equation}
Static mass is not a continuous parameter; it is a natural number.
There is no such thing as ``half an intermediate.''
The \emph{dynamic} mass incorporates the coupon-collector coverage-time
penalty: $M_{\mathrm{din}} = \kappa\,N\,H_N$, where
$H_N = \sum_{k=1}^{N} 1/k$ is the harmonic number
(see Demo~III, Step~6).

\medskip\noindent
\texttt{skyrmion.rs}: \texttt{mass\_gen1 = total / ps.len()}
where \texttt{total = ps.iter().map(|p|
p.intermediates.len()).sum()}.

\subsection*{L3. Law of 2-Hop Detection}
\phantomsection\label{law:2hop-detection}

For each core node $u$, traverse the path $u \to w \to v$
(forward-forward):
\[
  \text{belly}(u,v) = \text{children}(u) \cap \text{parents}(v).
\]
If $|\text{belly}| \geq 3$, a Prism
$\mathcal{P}(u, v, \text{belly})$ is committed. The optimal
partner is selected by belly size (greedy).

\medskip\noindent
\texttt{skyrmion.rs}: the main loop iterates over
\texttt{core\_nodes}, executes the 2-hop search with
\texttt{cands}, and computes the intersection via
\texttt{count\_isect(nb\_u, pv)} over sorted CSR in $O(D)$.

\subsection*{L4. Law of $K_5$ Contraction}
\phantomsection\label{law:k5-contraction}

If an external node $t$ satisfies the threat criterion
(Theorem~\ref{thm:kuratowski-threat}), it is absorbed into the
highest-degree pole:
\[
  \texttt{merge\_into}[t] = \arg\max\bigl(\deg(u),\deg(v)\bigr).
\]
Transitive merger chains are resolved by
integer pointer chasing:
\texttt{while merge\_into[t] != t \{ t = merge\_into[t]; \}}.

\subsection*{L5. Law of Generational Classification}
\phantomsection\label{law:generational-classification}

Prisms are classified by their \textbf{phase complexity}
(Definition~\texttt{def:topological-generation} from Volume~I).
For each intermediate $w_i$ of the Prism $\mathcal{P}$, the
\emph{causal phase} is the sign of its \emph{bulk momentum}:
\[
  \varphi(w_i) = \operatorname{sgn}\bigl(f^+(w_i) - f^-(w_i)\bigr)
  \;\in\; \{-1,\; 0,\; +1\}.
\]
The three \emph{phase classes} $\Phi^{-}, \Phi^{0}, \Phi^{+}$
collect the intermediates with negative, zero, and positive phase
respectively.  The generation is the number of occupied phase
classes:
\begin{equation}
  g(\mathcal{P}) = \bigl|\{\sigma \in \{-,0,+\} :
    \Phi^{\sigma}(\mathcal{P}) \neq \varnothing\}\bigr|
  \;\in\; \{1,\; 2,\; 3\}.
  \label{eq:law-generation}
\end{equation}
Generation~1 ($g=1$): all intermediates in a single phase class.
Generation~2 ($g=2$): two classes occupied.  Generation~3
($g=3$): all three classes present.  Three is the maximum possible
because $\varphi$ admits only three sign values.

\medskip\noindent
\texttt{skyrmion.rs}: \texttt{classify\_prism()} implements the phase-class
count over intermediates exclusively (poles excluded).
The function \texttt{prism\_signature()} (quartile signature) is
retained solely for diagnostic logging.

\subsection*{L6. Law of CPT Conjugation}
\phantomsection\label{law:cpt-conjugation}

CPT conjugation acts on a Prism by swapping its poles
$u \leftrightarrow v$, which reverses the direction of all
internal edges and, consequently, flips the sign of each causal
phase: $\varphi(w_i) \to -\varphi(w_i)$.  The \textbf{net causal
flux} (Definition~\texttt{def:matter-antimatter} from Volume~II)
\[
  \Phi(\mathcal{P}) = \sum_{i=1}^{N} \varphi(w_i)
\]
transforms as $\Phi \to -\Phi$.  A Prism is \textbf{matter}
if $\Phi > 0$ and \textbf{antimatter} if $\Phi < 0$.
Anti-Generation~1 comprises all Prisms with $g = 1$ and
$\Phi < 0$:
\begin{equation}
  \text{Anti-Gen}_1 = \bigl\{\mathcal{P} : g(\mathcal{P}) = 1
    \;\wedge\; \Phi(\mathcal{P}) < 0 \bigr\}.
  \label{eq:law-cpt}
\end{equation}

\medskip\noindent
\texttt{skyrmion.rs}: \texttt{classify\_prism()} computes $g$ and
$\Phi$ simultaneously; the matter/antimatter assignment
depends on the sign of $\Phi$.  The function
\texttt{anti\_signature()} is retained for diagnostic logging.

\stoneref{CPT conjugation is the logical contrapositive
(Theorem~\ref{thm:antimatter-tollens}).}

\subsection*{L7. Law of the Spectral Probe}
\phantomsection\label{law:spectral-probe}

The geometry of the graph is measured by random walkers
on the undirected Hasse diagram. A lazy walker
(\emph{lazy walk}) stays put with probability $1/2$ or jumps to a
uniform neighbor with probability $1/2$:
\begin{equation}
  \Pr[u_{t+1} = v \mid u_t = u] =
  \begin{cases}
    \frac{1}{2} & \text{if } v = u,\\[4pt]
    \frac{1}{2\deg(u)} & \text{if } (u,v) \in E.
  \end{cases}
  \label{eq:law-lazy}
\end{equation}
One measures $P(t) = W^{-1}\sum_i \delta_{u_t^{(i)}, u_0^{(i)}}$ and
$\dS(t) = -2\,d(\ln P)/d(\ln t)$.

\medskip\noindent
\texttt{spectral.rs}: \texttt{run\_walkers()} with
\texttt{rng.gen\_bool(0.5)}.

\subsection*{L8. Law of Causal Flux (Discrete Electromagnetism)}
\phantomsection\label{law:causal-flux}

\emph{Directed} walkers (cause $\to$ effect only) measure
transmission between topological signatures:
\begin{itemize}
  \item \textbf{Attraction}: Gen1 $\to$ AntiGen1 (opposite charges).
  \item \textbf{Repulsion}: Gen1 $\to$ Gen1 (like charges).
\end{itemize}
The asymmetry between both transmission probabilities is
the discrete analogue of the electromagnetic force.

\medskip\noindent
\texttt{spectral.rs}: \texttt{run\_transmission\_walkers()}
with mandatory movement (\texttt{pos = adj\_data[start +
rng.gen\_range(0..len)]}) and binary search over the
target sets.

% ═══════════════════════════════════════════════════════════════════════════════
\section{Theorems}
% ═══════════════════════════════════════════════════════════════════════════════

\begin{theorem}[Triangle-Freeness]
\label{thm:app-triangle-free}
The transitive reduction of any DAG is triangle-free.
If edges $u \to w$, $w \to v$, and $u \to v$ exist
simultaneously, then $u \to v$ is transitively redundant
and is eliminated. Therefore, $K_4$ cliques cannot exist
in the Hasse diagram, and $K_{2,N}$ (bipartite) is the maximal
permitted substructure.

\medskip\noindent
\texttt{skyrmion.rs}, line~35: ``Therefore $K_4$ cliques
(which require triangles) CANNOT exist---zero $K_4$ is not a bug,
it is a theorem.''
\end{theorem}

\begin{proof}
Let $H$ be the Hasse diagram (transitive reduction) of a DAG
$G = (V, \prec)$. Suppose, for contradiction, that there exist
$u, w, v$ with $u \covers w$, $w \covers v$, and
$u \covers v$ in $H$. But $u \prec w \prec v$ implies that
the edge $u \to v$ is transitively redundant (there exists a
2-hop path $u \to w \to v$), which contradicts its
membership in the transitive reduction. Therefore, $H$ does not
contain triangles. Since $K_4$ requires at least one
triangle, $K_4 \not\subset H$.
\end{proof}

\smallskip\noindent
{\small\itshape Note: this result is equivalent to the Topological
Non-Contradiction Principle
(Theorem~\ref{thm:non-contradiction},
Appendix~\ref{app:stone-duality}), which provides the dual
logical reading: the absence of triangles guarantees the consistency
of the deductive system.}

\begin{theorem}[2-Hop Completeness]
\label{thm:app-2hop}
The forward-forward search $u \to w \to v$ finds
\textbf{all} Causal Prisms $K_{2,N}$ in the core.
\end{theorem}

\begin{proof}
Let $u$ and $v$ be poles of a Prism with belly
$\{w_1,\ldots,w_N\}$. By definition, $u \covers w_i$ and
$w_i \covers v$ for all $i$. The path
$u \to w_i \to v$ guarantees that $v$ is discovered as a 2-hop
candidate through any $w_i$. The intersection
$\text{children}(u) \cap \text{parents}(v)$ recovers the complete
belly.

\medskip\noindent
\texttt{skyrmion.rs}: ``if $u$ and $v$ are poles of a timelike
$K_{2,N}$, then every belly node $w_i$ satisfies
$u \to w_i \to v$. The path guarantees $v$ is discovered---it is
exact.''
\end{proof}

\begin{theorem}[Computational Locality $O(N)$]
\label{thm:app-on}
Causal Prism detection scales linearly with the
number of events:
\begin{equation}
  T_{\text{detect}}
  = O\bigl(|C| \times \mathcal{V}_{\max}^2\bigr)
  = O(N),
  \label{eq:thm-on}
\end{equation}
where $|C| = N/10$ is the core (top $10\%$ by degree) and
$\mathcal{V}_{\max} = 15$ is the valence bound.
\end{theorem}

\begin{proof}
For each core node $u$, the 2-hop candidate collection
performs $O(D^2)$ operations where $D = \deg(u) \leq 15$. The
verification of each candidate (intersection of sorted
lists) is $O(D)$. The number of candidates per node is
$O(D^2) \leq 225$. Thus, the total work per node is
$O(D^2 \times D) = O(D^3)$, which is $O(1)$ since
$D \leq 15$. Summing over $|C|$ nodes:
$T = O(|C|) = O(N/10) = O(N)$.

\medskip\noindent
\texttt{skyrmion.rs}: ``$O(N)$ scaling IS the physics.
An $O(N^2)$ algorithm would imply non-local interactions---a
universe that could not scale.''
\end{proof}

\begin{theorem}[Causal Resolution (\ref{thm:causal-resolution})]
\label{thm:app-resolution}
The number of walkers needed to resolve the geometry
down to depth $t_{\max}$ with tolerance~$\epsilon$ is:
\begin{equation}
  W \geq \frac{t_{\max}^{\,\dS/2}}{c\;\epsilon^2}.
\end{equation}
For $\dS = 4$: $W \propto t_{\max}^2$ (quadratic cost).
See the complete proof in
Section~\ref{thm:causal-resolution}.

\medskip\noindent
\texttt{main.rs}: the number of walkers is fixed as
$W \sim N^2$ to guarantee constant signal-to-noise across the entire
diffusion range.
\end{theorem}

\begin{theorem}[Emergent Dimension $\dS \to 4$]
\label{thm:app-emergent-dim}
In the high-density limit ($N \gg 1$), the spectral
dimension of the causal set with Prisms converges to
$\dS = 4$ at intermediate diffusion times.
\end{theorem}

\begin{proof}[Sketch]
For a Poisson \emph{sprinkling} in a 4-dimensional
spacetime, the return probability satisfies
$P(t) \sim t^{-2}$~\cite{eichhorn2014}, whence $\dS = -2
\times (-2) = 4$. Causal Prisms stabilize $\dS$
by increasing $P(t)$ locally (the walker gets ``trapped''
in the intermediates), reducing the fluctuations of $\dS$
relative to the pure vacuum. The convergence of $\langle \dS \rangle$
to the value $4$ is verified numerically in the simulation
by averaging over $M$ independent realizations.
\end{proof}

\begin{theorem}[Poisson Volume (\ref{thm:volume-convergence})]
\label{thm:app-poisson-vol}
$\mathbb{E}[N] = \rho \cdot \mathrm{Vol}\bigl(\diam{A}{B}\bigr)$,
$\mathrm{Var}[N] = \mathbb{E}[N]$, $\sigma/\mu = 1/\sqrt{N}
\to 0$. The topological integral converges to the volume
integral of General Relativity~\cite{bombelli1987}.
\end{theorem}

\begin{theorem}[Myrheim--Meyer (\ref{thm:myrheim-meyer})]
\label{thm:app-mm}
The maximum chain length
$\ell(A,B)/\rho^{1/d} \to m_d\,\tau(A,B)$ as
$\rho \to \infty$~\cite{myrheim1978,brightwell1991}. The
covering-link count converges to the geodesic proper
time.
\end{theorem}

% ═══════════════════════════════════════════════════════════════════════════════
\section{Derived Results and Open Problems}
% ═══════════════════════════════════════════════════════════════════════════════

\subsection*{T1. Minimum Belly $N \geq 3$ (Derived)}
\phantomsection\label{thm-ref:min-belly}

The minimum belly of a Causal Prism is $N \geq 3$. This is
derived from the \textbf{Bifurcation-Convergence Uniqueness Theorem}
(Vol~II, Prism Uniqueness Theorem): in a $K_3$-free graph,
every structure with at least three independent $u$--$v$ paths of
exact length~2 is isomorphic to $K_{2,N}$ with $N \geq 3$. The bound is
tight: $K_{2,2}$ is a $C_4$ whose $S_2$ symmetry only generates
$U(1)$ (abelian group), insufficient for color confinement.
$S_3$ (the Weyl group of $SU(3)$) is required, which demands $N \geq 3$
(Vol~II, Unique Interaction Principle).

\medskip\noindent
\textit{Code}: \texttt{skyrmion.rs:73} ---
\texttt{const MIN\_PRISM\_SHARED: usize = 3}.

\subsection*{P1. Quantitative Mass Correspondence (Resolved)}
\phantomsection\label{open:mass-ratios}

The mass ratios $m_\mu/m_e \approx N_2/N_1$ and
$m_\tau/m_e \approx N_3/N_1$ should match
observations within $O(1)$ corrections arising from the
\emph{sprinkling} geometry.

\medskip\noindent
\textit{Status}: \textbf{resolved} by the occupation model
(\S\ref{sec:mass-hierarchy}). The phase fractions
$(p_+, p_0, p_-)$, derived analytically from the Poisson
\emph{sprinkling} geometry (Theorem~\ref{thm:phase-fractions}),
and the belly-size distribution $f(N)$
(Proposition~\ref{prop:belly-exponential}) completely determine
the mass ratios. The hierarchy $m_1 < m_2 < m_3$ is
structural for any $(p_+, p_0, p_-)$ with $p_i > 0$ and
$p_0 \ll p_\pm$. Numerically: $m_2/m_1 = 1{.}409$
(observed~$1{.}434$, error $1{.}8\%$) and $m_3/m_1 = 1{.}674$
(observed~$1{.}698$, error $1{.}4\%$). No parameter is fitted.

\subsection*{C2. Gravity as Connectivity Entropy}
\phantomsection\label{conj:gravity-entropy}

Einstein's equations $G_{\mu\nu} = 8\pi G\,T_{\mu\nu}$
emerge because the Prism density modifies the return
probability $P(t,u)$, and via the heat-kernel
expansion~\cite{minakshisundaram1949}, this modification maps to
the Ricci curvature. Newton's constant $G$ is identified
with a combination of the fundamental density~$\rho$ and the
topological inertia quantum~$\kappa$.

\medskip\noindent
\textit{Status}: Section~6.2 sketches the proof.
A proof sketch is presented in
Theorem~\ref{thm:gravity-inference-entropy}
(Appendix~\ref{app:stone-duality}, Nugget~IV) via
inferential entropy. The
numerical validation requires extracting $R_{\mu\nu}$ from $P(t,u)$
order by order in the heat expansion, which demands
$N > 10^7$ events to achieve sufficient precision.

\subsection*{D1. Unit-Charge Normalization (Definition)}
\phantomsection\label{def-ref:flux-normalization}

The normalization $f = F / |\mathcal{T}|$ isolates the intrinsic
coupling per target node. It is the discrete analogue of the
electric field $\mathbf{E} = \mathbf{F}/q$: the total force divided
by the test charge yields the field intensity per unit
charge, independent of the number of test charges.

The raw flux $\texttt{flux\_attraction}$ is dominated by the
combinatorial factor $|\text{Gen1}| \gg |\text{AntiGen1}|$. Dividing by
$|\text{targets}|$ removes this bias and reveals the fundamental
topological coupling (Vol~II, Definition: Flux per Charge).

\medskip\noindent
\textit{Code}: \texttt{spectral.rs:71--74},
\texttt{main.rs:259--262}.

\subsection*{P2. Coulomb Scaling $1/r^2$ (Open)}
\phantomsection\label{open:coulomb-scaling}

The transmission probabilities should reproduce the Coulomb
law ($1/r^2$) in the continuum limit. The
attraction/repulsion asymmetry is a topological CPT property.

\medskip\noindent
\textit{Status}: simulations consistently show the
expected asymmetry, but the $1/r^2$ scaling has not yet been
verified.

\subsection*{T2. Threat Threshold $K_{\text{threat}} = 2$ (Derived)}
\phantomsection\label{thm-ref:k-threat}

The value $K_{\text{threat}} = 2$ is derived from \textbf{Kuratowski's
Theorem}: a node connected to both poles and to $\geq 2$ intermediates
of a Prism $K_{2,N}$ forms a set of $2 + 2 + 1 = 5$ mutually
reachable vertices, threatening to produce a $K_5$ minor. The
absorption is the minimal mechanism that neutralizes this threat
(Theorem~\ref{thm:kuratowski-threshold}, \S1.4).

\medskip\noindent
\textit{Code}: \texttt{skyrmion.rs:77} ---
\texttt{const PRISM\_THREAT: usize = 2}.

\subsection*{P3. Confining Flux-Tube Area (Open)}
\phantomsection\label{open:flux-tube-area}

The confining flux-tube area must scale as the distance
between the poles in the emergent metric, in analogy with the QCD
flux tube.

\medskip\noindent
\textit{Status}: the absorbed-threat count
(\texttt{merge\_count}) is consistent with the prediction, but the
analytical proof of the area $\sim$ distance proportionality
remains open.

\subsection*{P4. Generational Exclusion ($g_{\max} = 3$) (Open)}
\phantomsection\label{open:gen-exclusion}

The sign function $\varphi \colon w_i \mapsto \{-1,\,0,\,+1\}$ has
exactly three output values. The generation number of a
prism $P = K_{2,N}$ is defined as $g(P) = |\{\text{distinct
values of } \varphi(w_i)\}| \leq 3$. The discrete topology of the Causal Prism
forbids a fourth generation: there is no fourth value of the
binary sign function that assigns chirality to the intermediates.

\medskip\noindent
\textit{Prediction}: no accelerator at any energy
will produce a fourth generational flavor.

\subsection*{P5. Absence of Primordial Gravitational Waves
  ($r < 10^{-3}$) (Open)}
\phantomsection\label{open:no-tensor}

In the UV regime, $\dS \to 2$ (Spectral Convergence Theorem,
\S3.1). In two dimensions the Weyl tensor vanishes identically:
there are no propagating gravitational degrees of freedom. Therefore, the
Kuratowski epoch (pre-topological inflation) does not generate
tensor perturbations.

\medskip\noindent
\textit{Prediction}: the CMB tensor-to-scalar ratio satisfies
$r < 10^{-3}$, beyond the reach of LiteBIRD if topological
inflation dominates.

\subsection*{P6. The $\alpha$--$\Omega$ Lock}
\phantomsection\label{open:alpha-omega-lock}

We define the \emph{dark self-energy ratio}:
\[
  \Omega_{\mathrm{energy}}
  \;=\; \frac{\sum N_P^{\,2} - \sum|\Phi(P)|^2}{\sum|\Phi(P)|^2}
  \;=\; \frac{1}{\mathcal{Q}_{\mathrm{topo}}} - 1\,,
\]
where $\mathcal{Q}_{\mathrm{topo}} = \sum|\Phi|^2/\sum N^2$ is the
fraction of gravitational self-energy that is electromagnetically
active.  Since $\alpha_0 = \mathcal{Q}_{\mathrm{topo}}/(8\pi)$
(C6), one obtains:
\[
  \boxed{\;\alpha_0\,(1 + \Omega_{\mathrm{energy}}) = \frac{1}{8\pi}\;\,.}
\]

\medskip\noindent
\textit{Nature of the identity.}\enspace
This relation is an \textbf{algebraic identity}: both sides are
functions of the same quantity $\mathcal{Q}_{\mathrm{topo}}$, so
it is satisfied by construction at each~$N$. Its non-trivial physical
content is that a single topological quantity---the fraction
$\mathcal{Q}_{\mathrm{topo}}$---simultaneously parametrizes the
electromagnetic coupling and the dark energy fraction.
The physical prediction resides in the \textbf{value} of
$\mathcal{Q}_{\mathrm{topo}}$ in the limit $N \to \infty$.

\medskip\noindent
\textit{Analytical result}
(\S\ref{sec:em-coupling}): the port count of $K_{2,3}$
establishes $\mathcal{Q}_{\mathrm{topo}} = 4/21$ for an isolated Prism.
Hence,
\[
  \alpha_0 = \frac{1}{42\pi} \approx \frac{1}{131{.}95}
  \qquad\text{(bare coupling)},
  \qquad
  \Omega_{\mathrm{energy}} = \frac{21}{4} - 1 = \frac{17}{4} = 4{.}25\,.
\]
The laboratory-measured fine-structure constant
($\alpha^{-1} = 137{.}036$) is weaker than $\alpha_0$ due to
vacuum polarization: virtual pairs $K_{2,3} +
\bar{K}_{2,3}$ and the baryonic sector $K_{3,3}$ screen the bare
charge at long distances. The residual $\alpha_0 \to \alpha_{\text{lab}}$
($\Delta\alpha^{-1} \approx 5$) quantifies the contribution of these
sectors.

\medskip\noindent
\textit{Finite-size scaling.}\enspace
The collective value $\mathcal{Q}_{\mathrm{topo}}(N)$ differs from $4/21$
by planarity correlations (Theorem~\ref{thm:planarity-screening}).
Finite-size scaling (FSS) extrapolates
$\Omega_\infty = 5{.}57$ versus the Planck~2018 value
of~$5{.}36$---an agreement of $\sim\!4\%$ that improves with~$N$.

\medskip\noindent
\textit{Prediction}: if the identity reflects a fundamental law,
any measurement of $\alpha$ as a function of redshift
$z$ must satisfy $\alpha(z)\,[1 + \Omega(z)] = 1/(8\pi)$.

\subsection*{P7. Proton Stability
  ($\tau_p > 10^{40}$~years) (Open)}
\phantomsection\label{open:proton-stability}

The baryon number is the linking number
of the prism knot in the Hasse diagram. Proton decay
requires simultaneously unknotting $\sim\!10^{60}$ causal
links. The probability of such a topological fluctuation is
$\sim e^{-N}$, where $N$ is the number of intermediates
involved.

\medskip\noindent
\textit{Prediction}: $\tau_p > 10^{40}$~years, consistent with
the experimental bounds from Super-Kamiokande and future
Hyper-K detectors.

\subsection*{C5. Decoherence Threshold}
\phantomsection\label{conj:decoherence-threshold}

When the number of internal causal paths of a subsystem
$\mathcal{S}$ violates the bound $W_{\mathcal{S}} < t^2 /
(c\,\epsilon_{\text{coh}}^2)$, the local spectral dimension
ceases to converge to~4 and the continuum description breaks down.
Quantum decoherence~\cite{zurek2003} is the loss of
causal resolution (Theorem~\ref{thm:decoherence-resolution}).

\medskip\noindent
\textit{Status}: \textbf{proven} in Theorem~\ref{thm:decoherence-resolution} (Epilogue).
Retained here as a historical reference.

\subsection*{C6. Dark Matter as Phase Decomposition (Theorem)}
\phantomsection\label{conj:dark-matter}

The gravitational self-energy of each prism $P = K_{2,N}$ is
$N^2$ (pairwise interactions of the $N$ intermediates); its
electromagnetic self-energy is $|\Phi(P)|^2$ (only the net
coherent phase contributes to the directed flux). The
topological charge fraction is
$\mathcal{Q}_{\mathrm{topo}} = \sum|\Phi(P)|^2 / \sum N_P^{\,2}$.

The bare fine-structure constant emerges as
$\alpha_0 = \mathcal{Q}_{\mathrm{topo}} / (8\pi)$, where the factor
$8\pi$ comes from the isotropic projection of the double light cone
(Vol~II, Appendix~B). The port count of $K_{2,3}$
(\S\ref{sec:em-coupling}) establishes
$\mathcal{Q}_{\mathrm{topo}} = 4/21$ for an isolated Prism, whence
$\alpha_0 = 1/(42\pi)$.

The dark-to-visible matter ratio is obtained from the
\textbf{self-energy decomposition}: the dark sector is the
gravitational self-energy not accounted for by electromagnetic
interactions:
$\Omega_{\mathrm{energy}} = (\sum N^2 - \sum|\Phi|^2) / \sum|\Phi|^2
= 1/\mathcal{Q}_{\mathrm{topo}} - 1$.
The algebraic identity $\alpha_0\,(1 + \Omega) = 1/(8\pi)$
(\S\ref{open:alpha-omega-lock}) links both quantities through
$\mathcal{Q}_{\mathrm{topo}}$.

For an isolated Prism ($\mathcal{Q}_{\mathrm{topo}} = 4/21$):
$\alpha_0^{-1} = 42\pi \approx 131{.}95$ (bare coupling) and
$\Omega_{\mathrm{energy}} = 17/4 = 4{.}25$. The Planck~2018
observation is $\Omega \approx 5{.}36$, a discrepancy of
${\sim}\!21\%$. The residual quantifies the dynamic mass
contribution from baryonic confinement ($K_{3,3}$) and collective
planarity effects; its first-order calculation remains
open.

\medskip\noindent
\textit{Status}: \textbf{proven} (Phase Coherence
and Emergent Fine-Structure Constant Theorems, Vol~II).

% ═══════════════════════════════════════════════════════════════════════════════
\section{Architecture $\leftrightarrow$ Physics Correspondence}
% ═══════════════════════════════════════════════════════════════════════════════

\noindent
The following table summarizes how each component of the simulation
engine implements a fundamental physical operation:

\medskip
\begin{center}
\small
\begin{tabular}{p{4.5cm}ll}
\toprule
\textbf{Module / Function} & \textbf{Operation} &
  \textbf{Physics} \\
\midrule
\texttt{diamond::sprinkle()}
  & Poisson sprinkling $\rho = 1$
  & Spacetime genesis \\
\texttt{diamond::is\_causal()}
  & $dt^2 > |\Delta\vec{x}|^2$
  & Lorentzian causality \\
\texttt{diamond::build\_hasse\_*()}
  & Transitive reduction
  & $dx \to \covers$ \\
\texttt{skyrmion::apply\_defect()}
  & $K_{2,N}$ detection
  & Matter creation \\
\texttt{skyrmion} (threat)
  & $K_5$ contraction
  & Strong nuclear force \\
\texttt{skyrmion::prism\_signature()}
  & Signature $[q_0,q_1,q_2,q_3]$
  & Leptonic flavor \\
\texttt{skyrmion::anti\_signature()}
  & $\bar{\sigma}_i = -\sigma_{3-i}$
  & Antimatter (CPT) \\
\texttt{spectral::run\_walkers()}
  & $P(t), \dS(t)$
  & Spacetime dimension \\
\texttt{spectral::run\_transmission\_*()}
  & Directed flux
  & Electromagnetism \\
\texttt{main::run\_realization()}
  & $M$ parallel ensembles
  & Everett universes \\
\texttt{main::average\_ensemble()}
  & $\langle P(t)\rangle$
  & Thermodynamic limit \\
\bottomrule
\end{tabular}
\end{center}

\bigskip\noindent
\begin{result}[title={Logical Efficiency as the Engine of
Physics}]
\phantomsection\label{res:logical-efficiency}
In classical physics, the engine that selects the physical
trajectory is the principle of least action: $\delta S = 0$.
In the Kuratowski Calculus, this principle is replaced by
\textbf{Logical Efficiency}---the avoidance of topological
contradictions. Every operation of the engine (\texttt{Phases 1--3})
reduces to a single imperative: \emph{maintain the consistency
of the DAG}. The transitive reduction eliminates redundancies (Law~L1).
Prism detection identifies the maximal structure
compatible with triangle-freeness
(Theorem~\ref{thm:app-triangle-free}). $K_5$ contraction
preserves planarity (Theorem~\ref{thm:kuratowski-threat}). The
$O(N)$ scaling demonstrates that this consistency is maintainable
without non-local information
(Theorem~\ref{thm:app-on}). The classical action $S[\gamma]
= -|\gamma|$ (Definition~\ref{def:discrete-action}) is the
variational consequence of this imperative: the geodesic is the
chain that maximizes causal flux without violating the topology. Physics
does not obey laws imposed from outside; it obeys the
only law that a finite, causal, planar graph can obey:
its own logical consistency.
\end{result}

% ███████████████████████████████████████████████████████████████████████████████
%   APPENDIX B: STONE DUALITY — PHYSICS AS LOGIC
% ███████████████████████████████████████████████████████████████████████████████

\chapter{Stone Duality: Physics as Logic}
\label{app:stone-duality}

\begin{introduction}[An Unexpected Lens]
  \item This appendix adds no new physics to the Kuratowski
        Calculus; it reveals a \textbf{deep logical structure} that
        was already present in every axiom, law, and theorem of the
        preceding chapters.
  \item Stone Duality~\cite{stone1936,stone1937,johnstone1982}
        establishes a categorical equivalence between topologies and
        Boolean algebras (propositional logic). If
        spacetime is a topology (Alexandrov topology on a DAG),
        then spacetime \textbf{is} a logical system.
  \item We present six results---\emph{Golden Nuggets}---that
        emerge from reading the formalism through this duality.
\end{introduction}

\noindent
The fundamental correspondence articulated in this appendix is:

\medskip
\begin{center}
\begin{tabular}{p{4.5cm}ll}
\toprule
\textbf{Topology (Kuratowski)} & \textbf{Logic (Stone)} &
  \textbf{Physics}\\
\midrule
Event $e \in V$ (Def.~\ref{def:event})
  & Proposition $\varphi$ & Spacetime point\\
Covering link $\covers$ (Def.~\ref{def:covering})
  & Implication $\vdash$
  & Causal differential $dx$\\
Diamond $\diam{A}{B}$ (Def.~\ref{def:alexandrov})
  & Deductive theory $\Gamma$
  & Integration region\\
Prism $K_{2,N}$ (Thm.~\ref{thm:bipartite-matter})
  & Auxiliary lemma ($N$ steps)
  & Particle (mass $N$)\\
Threat $K_5$ (Thm.~\ref{thm:kuratowski-threat})
  & Logical paradox
  & Topological singularity\\
$K_5$ contraction (Law~\ref{law:k5-contraction})
  & Resolution of the paradox
  & Strong nuclear force\\
Maximal chain $\gamma$ (Def.~\ref{def:geodesic})
  & Optimal proof
  & Geodesic (proper time)\\
Antichain $\mathcal{A}$ (Def.~\ref{def:antichain})
  & Independent axioms
  & Cauchy hypersurface\\
\bottomrule
\end{tabular}
\end{center}

% ═══════════════════════════════════════════════════════════════════════════════
\section{Nugget I: The Stone--Kuratowski Bridge}
% ═══════════════════════════════════════════════════════════════════════════════

\begin{theorem}[Stone--Kuratowski Duality]
\label{thm:stone-kuratowski}
Let $\mathcal{C} = (V, \prec)$ be a finite causal set. The
Alexandrov topology~\cite{sorkin1991topology} on $\mathcal{C}$
(whose open sets are the \emph{past-closed} sets:
$x \in U \wedge y \prec x \Rightarrow y \in U$) is exactly the
spectral topology of the free Heyting algebra generated by~$V$.
In particular:

\begin{enumerate}
  \item The covering link $u \covers v$ \textbf{is} the
    logical implication $\varphi_u \vdash \varphi_v$: ``the truth
    of event~$u$ causally implies the truth of event~$v$.''
  \item The Alexandrov diamond $\diam{A}{B} = \{x : A \prec x
    \prec B\}$ \textbf{is} the deductive theory closed under
    implication between propositions $\varphi_A$ and~$\varphi_B$.
  \item The transitive reduction (Hasse diagram,
    Theorem~\ref{thm:covering-differential}) \textbf{is}
    proof normalization (cut
    elimination)~\cite{gentzen1935}: it eliminates every implication
    $\varphi_u \vdash \varphi_v$ that is derivable by
    transitivity, retaining only the atomic inferential
    steps.
\end{enumerate}
\end{theorem}

\begin{proof}
The Alexandrov topology on a finite poset $P$ assigns to each
$x \in P$ the basic open set $\downarrow\!x = \{y : y \preceq x\}$.
The function $x \mapsto \downarrow\!x$ is an order isomorphism
between $(P, \preceq)$ and the lattice of open sets $(\mathcal{O}(P),
\subseteq)$. By Stone's theorem for Heyting
algebras~\cite{stone1937,vickers1989}, this lattice is the
Lindenbaum--Tarski algebra of an intuitionistic logic whose
atomic propositions are the elements of~$V$.

For part (3), observe that in intuitionistic logic, Gentzen's
\emph{Hauptsatz}~\cite{gentzen1935} guarantees that every
proof can be normalized by eliminating ``cuts'' (uses of
transitivity). The transitive reduction $u \covers v$ retains
$u \to v$ if and only if no $w$ exists such that $u \prec w \prec v$:
that is, it retains exactly the \emph{atomic} implications, which
are the already-normalized proofs. The elimination of transitive
edges in the DAG is literally cut elimination in the
sequent calculus of causal logic.
\end{proof}

\begin{result}[title={Spacetime Is a Deductive System}]
\phantomsection\label{res:stone-deductive-system}
The Kuratowski Calculus is not an \emph{analogy} with logic:
it \textbf{is} logic. Each region of spacetime is a deductive
theory closed under causal implication. The physics we measure is
the semantics (the models) of that formal system. Stone
Duality~\cite{stone1936} guarantees that both descriptions are
categorically equivalent: to know the topology is to know the
logic, and vice versa.
\end{result}

% ═══════════════════════════════════════════════════════════════════════════════
\section{Nugget II: Non-Contradiction = Triangle-Freeness}
% ═══════════════════════════════════════════════════════════════════════════════

\begin{theorem}[Topological Non-Contradiction Principle]
\label{thm:non-contradiction}
Let $H = (V, E_H)$ be the Hasse diagram of a causal set
$\mathcal{C}$. Then $H$ is free of directed triangles.
This property is the topological realization of the \textbf{Principle
of Non-Contradiction}: in the dual logical system, no
proposition can be simultaneously an axiom and a derived theorem
with respect to another proposition.
\end{theorem}

\begin{proof}
Suppose for contradiction that a directed triangle
$u \covers w \covers v$ with $u \covers v$ also
present in~$H$ exists. But $u \covers v$ is a covering link,
which by definition requires that no $w$ with
$u \prec w \prec v$ exists. The presence of $u \covers w \covers v$
violates this definition. Contradiction. Therefore $H$ contains no
directed triangles. $\square$

\medskip\noindent
\textit{Logical reading.} In the dual system, $u \covers v$
means ``$\varphi_u$ implies $\varphi_v$ in a single
inferential step (implication axiom).'' If a path
$\varphi_u \vdash \varphi_w \vdash \varphi_v$ existed AND simultaneously a
direct axiom $\varphi_u \vdash \varphi_v$, then the same
logical link would be simultaneously primitive (axiomatic) and
derived (via~$w$). This is a logical contradiction at the
meta-level: one and the same judgment cannot be both assumed and
proven. The transitive reduction structurally forbids
it---implementing the logical coherence of the system.

\smallskip\noindent
\textit{Note:} the purely combinatorial proof of triangle-freeness
is established independently in
Theorem~\ref{thm:app-triangle-free} (Appendix~\ref{app:axioms}).
The present result adds the \emph{logical reading}: the same
graph-theoretic property is the Principle of Non-Contradiction of the
dual deductive system.
\end{proof}

\begin{result}[title={The Geometry of Spacetime Obeys Classical
Logic}]
\phantomsection\label{res:stone-non-contradiction}
The absence of triangles in the Hasse diagram is not a
combinatorial curiosity: it is the guarantee that spacetime,
understood as a deductive system, is \textbf{logically
consistent}. The transitive reduction
(Theorem~\ref{thm:covering-differential}) fulfills in the causal DAG the
same function that the Principle of Non-Contradiction fulfills in
formal logic: it prevents a single relation from being
simultaneously atomic and derived. The simulation engine
(\texttt{build\_hasse\_sparse}) executes this prohibition as its
most expensive operation ($O(N^2)$ for $N \leq 15\,000$,
$O(N^{1.5})$ for $N > 15\,000$)---logical consistency is the
most expensive computation the universe performs.
\end{result}

% ═══════════════════════════════════════════════════════════════════════════════
\section{Nugget III: Antimatter as \emph{Modus Tollens}}
% ═══════════════════════════════════════════════════════════════════════════════

\begin{theorem}[Matter--Antimatter Duality as the
Logical Contrapositive]
\label{thm:antimatter-tollens}
Let $\Pi = (u, v, \{w_1, \ldots, w_N\})$ be a Causal Prism with
topological signature $\sigma = [q_0, q_1, q_2, q_3]$. The CPT
conjugation
\begin{equation}
  \bar{\sigma}_i = -\sigma_{3-i} + 16
  \label{eq:cpt-contrapositive}
\end{equation}
(\texttt{skyrmion::anti\_signature()}) is the \textbf{logical
contrapositive} of the Prism's inferential chain. Specifically:

\begin{enumerate}
  \item The Prism $\Pi$ encodes the \emph{Modus Ponens} chain:
    \begin{equation}
      \varphi_u \;\vdash\; \varphi_{w_1}, \ldots, \varphi_{w_N}
      \;\vdash\; \varphi_v.
    \end{equation}
    ``From premise $u$, through $N$ intermediate lemmas, one
    concludes~$v$.''
  \item The anti-Prism $\bar{\Pi}$ encodes the \emph{Modus Tollens}
    (contrapositive):
    \begin{equation}
      \neg\varphi_v \;\vdash\; \neg\varphi_{w_N}, \ldots,
      \neg\varphi_{w_1} \;\vdash\; \neg\varphi_u.
    \end{equation}
    ``If conclusion $v$ is false, then all
    intermediates are false and premise $u$ is false.''
  \item \textbf{Annihilation} $\Pi + \bar{\Pi} \to \text{vacuum}$
    is tautological collapse: the inferential chain meets
    its contrapositive, the $N$ intermediates cancel, and
    the released energy $E \propto N$ is the \textbf{combinatorial
    resolution} of a tautology---the proof of
    $(\varphi \Rightarrow \psi) \Leftrightarrow
    (\neg\psi \Rightarrow \neg\varphi)$.
\end{enumerate}
\end{theorem}

\begin{proof}
The contrapositive is the operation $(\varphi \Rightarrow \psi)
\mapsto (\neg\psi \Rightarrow \neg\varphi)$, valid in
both classical and intuitionistic logic. In the dual poset, negating a
proposition amounts to taking the complement of its open filter,
which reverses the order and flips the sign of the quartile
position. The signature $\sigma = [q_0, q_1, q_2, q_3]$ encodes the
statistical distribution of the intermediates across four quartiles
of the causal bulk. The contrapositive reverses the temporal order
($i \mapsto 3 - i$) and reflects the position within each quartile
($q \mapsto -q + 16$, where $16$ is the \texttt{u32} packing
modulus). This is precisely~\eqref{eq:cpt-contrapositive}.

Annihilation releases the $N$ intermediates because the tautology
$(\varphi \Rightarrow \psi) \wedge (\neg\psi \Rightarrow
\neg\varphi)$ is valid without hypotheses---the proof ``collapses''
to zero steps, returning the $N$ covering links to the vacuum.
\end{proof}

% ═══════════════════════════════════════════════════════════════════════════════
\section{Nugget IV: Gravity as Inferential Entropy}
% ═══════════════════════════════════════════════════════════════════════════════

\begin{theorem}[Einstein's Equations as Inferential Thermodynamics]
\label{thm:gravity-inference-entropy}
Let $\nu(x)$ be the local density of Causal Prisms around an
event~$x$, and let $\Omega(x)$ be the number of topologically
distinct causal chains (proofs) connecting the past
to the future of~$x$ through its Prisms. We define the
\textbf{local inferential entropy}:
\begin{equation}
  S_{\text{inf}}(x) \;=\; \ln \Omega(x).
  \label{eq:inferential-entropy}
\end{equation}
Then, in the continuum limit ($\rho \to \infty$), the gradient
of $S_{\text{inf}}$ reproduces Einstein's equations:
\begin{equation}
  G_{\mu\nu}(x) \;=\; 8\pi G \; T_{\mu\nu}(x)
  \;\;\longleftrightarrow\;\;
  \nabla^2 S_{\text{inf}} \;=\; 8\pi G \;\nu(x),
  \label{eq:einstein-as-entropy}
\end{equation}
where $T_{\mu\nu}$ is identified with the topological defect
density $\nu(x)$ (Conjecture~C2, Appendix~A).
\end{theorem}

\begin{proof}[Sketch]
The key is the relationship between the heat-kernel
expansion~\cite{minakshisundaram1949,gilkey1975} and the
random-walker return probability
(Definition~\ref{def:spectral-dimension}):
\[
  P(t, x) \;=\; \frac{1}{(4\pi t)^{d/2}}
  \Bigl(1 - \tfrac{1}{6}R(x)\,t + O(t^2)\Bigr),
\]
where $R(x)$ is the scalar curvature. A region with high Prism
density $\nu(x)$ traps walkers (the intermediates act as local
\emph{loops}), increasing $P(t,x)$ relative to the vacuum. This
increase corresponds to a negative contribution to the term
$-\tfrac{1}{6}R(x)t$, that is, to positive curvature. But the
number of ways to trap a walker is exactly
$\Omega(x)$: each Prism offers $N$ alternative paths through which
the walker can return. Thus, $P(t,x) \propto \Omega(x)$, and
therefore $R(x) \propto -\ln \Omega(x) = -S_{\text{inf}}(x)$.

Translated to the logical dual: a region with many Prisms has
more \textbf{possible proofs} connecting premises to
conclusions---greater inferential richness. This richness curves
spacetime. \textbf{Gravity is the tendency of the logical
system to maximize its inferential entropy}: events migrate
(in the \emph{sprinkling} sense) toward regions where there are more
ways to prove theorems.
\end{proof}

\begin{result}[title={Gravity as the Second Law of
Logic}]
\phantomsection\label{res:stone-gravity-logic}
Just as the second law of thermodynamics states that
physical systems evolve toward maximum molecular disorder,
gravity states that the logical system of spacetime
evolves toward the maximum number of possible proofs.
Mass curves space because it \textbf{adds lemmas} to the
deductive system: each intermediate of a Prism is a lemma that opens new
inferential pathways. The more lemmas, the more proofs;
the more proofs, the greater the inferential entropy; the greater
the inferential entropy, the greater the curvature. Einstein did not discover a
physical law: he discovered the \textbf{Second Law of
Causal Logic}.
\end{result}

% ═══════════════════════════════════════════════════════════════════════════════
\section{Nugget V: Black Holes as G\"odelian Horizons}
% ═══════════════════════════════════════════════════════════════════════════════

\begin{theorem}[Event Horizon as the G\"odel
Limit]
\label{thm:godel-horizon}
Let $\mathcal{B} \subset V$ be a region of the causal DAG such that its
Prism density satisfies the Bekenstein
bound~\cite{bekenstein1981}:
\begin{equation}
  \nu(\mathcal{B}) \cdot |\mathcal{B}|
  \;>\;
  \frac{|\partial \mathcal{B}|}{4\,\lP^2},
  \label{eq:bekenstein-logical}
\end{equation}
where $|\partial \mathcal{B}|$ is the number of covering links
crossing the boundary of~$\mathcal{B}$. Then,
$\mathcal{B}$ is a \textbf{G\"odelian domain}: it contains
propositions (events) whose truth cannot be proven
(causally reached) by any inferential chain originating
outside~$\mathcal{B}$.
\end{theorem}

\begin{proof}[Sketch]
Bound~\eqref{eq:bekenstein-logical} establishes that the number of
``internal theorems'' (events causally connected within
$\mathcal{B}$, each multiplied by the $N$ inferential pathways
of its Prisms) exceeds the capacity of the boundary to
transmit proofs. The $|\partial \mathcal{B}|$ boundary links
are the only channels through which an external inferential
chain can enter~$\mathcal{B}$. When the internal inferential
information exceeds this bandwidth, there necessarily exist
interior propositions $\varphi_x$ such that
$\varphi_x$ is \emph{true} (the event $x$ exists in the DAG)
but \emph{unprovable} from the exterior---the exact definition
of a G\"odelian sentence~\cite{godel1931}.

The event horizon is the surface $\partial \mathcal{B}$ where
the inferential bandwidth saturates. Outside the horizon, the
deductive system is complete (every reachable proposition is
provable). Inside, the inferential richness exceeds the
communication capacity: there are causal truths that no external observer
can derive.
\end{proof}

\begin{result}[title={The Information Paradox as
Incompleteness}]
\phantomsection\label{res:stone-information-paradox}
The so-called ``information paradox'' of black
holes~\cite{hawking1975} dissolves under Stone Duality. There
is no paradox because there was never any loss: the propositions interior
to the horizon are \emph{true but unprovable} from the
exterior---exactly like G\"odelian sentences in
arithmetic~\cite{godel1931}. Hawking radiation is the
analogue of the \emph{partial theorems} that the external system
can infer about the inaccessible region: it provides statistical
information (thermal entropy) but cannot reconstruct the
individual proofs (microstates). The Bekenstein limit
$S = A/(4\lP^2)$ counts exactly the number of interior
propositions that are \textbf{true-but-unprovable} as seen
from the boundary: it is the G\"odelian incompleteness index of the
region.
\end{result}

% ═══════════════════════════════════════════════════════════════════════════════
\section{Nugget VI: Parity Violation as Logical Holomorphy}
% ═══════════════════════════════════════════════════════════════════════════════

\begin{theorem}[Parity Violation as Holomorphic Consistency]
\label{thm:parity-stone}
Let $(P, \prec)$ be a finite causal DAG and let $\Pi = (u, v,
\{w_1, \ldots, w_N\})$ be a Prism with a Kuratowski twist at the future
pole~$v$. In Stone's logical reading:
\begin{enumerate}
  \item The causal order $\prec$ induces a \textbf{direction
    of the deductive system}: \emph{Modus Ponens} flows from premises
    to conclusions ($\varphi_u \vdash \varphi_v$). This direction
    is the logical analogue of the complex (holomorphic) structure
    of the Belyi surface.
  \item Both twists---left-handed and right-handed---produce valid dessins
    on orientable surfaces. There is no ``non-orientability''
    in Grothendieck's theory: every pair $(\sigma_0, \sigma_1)$
    generates an orientable surface.
  \item A \textbf{left-handed twist} corresponds to a
    \emph{valid non-constructive lemma}: a detour in the inferential
    chain whose complex structure aligns with the direction
    of deduction. The detour adds complexity (genus $> 0$)
    but is \textbf{holomorphic}---accessible to the
    simplification operator (the edge flip).
  \item A \textbf{right-handed twist} corresponds to a detour whose
    complex structure \textbf{opposes} the direction of
    deduction. The detour is \textbf{anti-holomorphic}: it requires
    complex conjugation ($z \to \bar{z}$) to map to the
    continuum. It exists as a valid topological state, but
    is \emph{invisible} to the simplification operator.
  \item The \textbf{holomorphic veto}
    (Theorem~\ref{thm:parity-veto}): Tarski's consequence
    operator~\cite{tarski1930} is strictly causal---it acts
    in the direction $\varphi \vdash \psi$, never in the opposite one.
    Topologically: the Belyi function is holomorphic.
    Logically: deductions flow only toward the future.
    Anti-holomorphic states are not \emph{forbidden}; they are
    \emph{inaccessible} to the deductive machinery.
  \item For \textbf{antimatter} (contrapositive,
    Theorem~\ref{thm:antimatter-tollens}): the direction of
    deduction is reversed ($\neg\varphi_v \vdash \neg\varphi_u$),
    and with it the complex structure is conjugated. Right-handed
    twists become holomorphic (accessible to \emph{Modus
    Tollens}); left-handed ones become anti-holomorphic (inaccessible).
\end{enumerate}
\end{theorem}

\begin{proof}[Sketch]
By Stone Duality, the open sets of the Alexandrov topology
on $(P, \prec)$ correspond to filters of the Boolean
algebra of causal propositions. The complex structure of the
Riemann surface (Belyi) translates, via this duality,
into the \textbf{directionality} of Tarski's consequence
operator~\cite{tarski1930}: if $\varphi \vdash \psi$, the
proof flows from $\varphi$ to $\psi$ (holomorphic), never in
the opposite direction (anti-holomorphic) within the same chain.

A left-handed twist produces a detour whose logical structure
respects the direction of the deductive flow: the
proof-simplification operator (the logical analogue of the weak force)
can act on it, eliminating the detour. A right-handed twist
produces a detour whose logical structure opposes the flow:
the simplification operator cannot access it---not because
it is inconsistent, but because its complex geometry is the conjugate
of the operator's geometry.

CPT conjugation reverses the direction of the consequence
operator ($\vdash$ is replaced by its contrapositive), conjugating
the complex structure and swapping the accessible chiralities.
\end{proof}

\begin{result}[title={Parity Violation Is Holomorphic Consistency}]
\phantomsection\label{res:stone-parity}
Parity violation is not an accidental asymmetry of the
universe: it is \textbf{holomorphy of the deduction operator}. The
deductive system (spacetime, read as a proof algebra)
admits detours of both chiralities as valid states, but its
simplification operator only accesses those that align with the
inferential arrow. For matter (\emph{Modus Ponens}), the
accessible chirality is left-handed. For antimatter (\emph{Modus
Tollens}), right-handed. The operator $P_L = \tfrac{1}{2}(1-\gamma^5)$
is not an empirical parameter of the Standard Model: it is the
\textbf{holomorphic projection} of the causal deductive system.
\end{result}

\bigskip
\begin{center}
\large\itshape\color{structure}
``Stone Duality reveals that the Kuratowski Calculus\\
does not \emph{model} reality with logic:\\
reality \textbf{is} logic executing itself.\\
Each particle is a theorem, each force a rule of inference,\\
each black hole a horizon of incompleteness,\\
each chiral asymmetry a consistency requirement.\\
The universe does not obey laws---it \emph{is} the law.''
\end{center}

% ███████████████████████████████████████████████████████████████████████████████
%   BIBLIOGRAPHY
% ███████████████████████████████████████████████████████████████████████████████
\backmatter
\addcontentsline{toc}{chapter}{Bibliography}
\bibliography{modulo_synthesis_kuratowski_calculus_en}

\end{document}
